\documentclass[11pt, a4paper]{article}

% --- SIMPLIFIED PREAMBLE ---
\usepackage[a4paper, top=2.5cm, bottom=2.5cm, left=2cm, right=2cm]{geometry}
\usepackage{graphicx}
\usepackage{amsmath, amssymb, amsthm}
\usepackage{amsfonts}
\usepackage{algorithm}
\usepackage{algpseudocode}
\usepackage{bm}
\usepackage{geometry}
\usepackage{hyperref}
\usepackage{tikz}
 
\newcommand{\R}{\mathbb{R}}
\newcommand{\layernorm}{\mathrm{LayerNorm}}
\newcommand{\softmax}{\mathrm{Softmax}}
\newcommand{\GRU}{\mathrm{GRU}}
\newcommand{\MLP}{\mathrm{MLP}}
\newcommand{\WM}{\mathrm{WeightedMean}}

\title{Equations}
\author{ Ricker }
\date{June 2026}

\newcommand{\D}{\mathrm{d}}

\setlength{\parindent}{0pt}

\begin{document}
\section{Neurons}

All state is fixed-point and updated \emph{lazily}: a compartment only advances when an event
touches it, decaying from its last-event timestamp $\tau$ to the current time $t$. Time is the
$\mathfrak{u}16$ event clock (ticks).

\subsection{Base-2 decay operator}
Every leak is the same $O(1)$ base-2 approximation of an exponential (\texttt{math::decay}):
\begin{equation}
    \mathcal{D}_k(v, \Delta t) = v \cdot 2^{-\Delta t / 2^{k}},
\end{equation}
i.e.\ $v$ halves every $2^{k}$ ticks (half-life $2^{k}$). The exponent $k$ selects the time constant
of each compartment; the same operator decays alpha, branch voltage, and the soma potential.

\subsection{Synapse}
Let $x \in \mathfrak{u}8$ (position along the dendrite), $w \in \mathfrak{i}8$ (weight),
$\alpha \in \mathfrak{u}8$ (activity / eligibility trace), $\tau \in \mathfrak{u}16$ (last-event time).
\begin{equation}
    s = (x,\,w,\,\alpha,\,\tau)
\end{equation}
Synapses on a dendrite are ordered by strictly increasing, unique $x$. The trace $\alpha$ decays with
$k_\alpha = 11$ and is boosted by $\alpha_b = 64$ each time a forward action potential arrives:
\begin{equation}
    \alpha \leftarrow \mathcal{D}_{k_\alpha}\!\big(\alpha,\, t - \tau\big), \qquad
    \alpha \leftarrow \min(\alpha + \alpha_b,\; 255).
\end{equation}

\subsection{Basal dendrite}
Synapses more distal than $i$ (larger $x$) gate $i$'s contribution: their (decayed) traces sum, with
weight that falls off over $x$ with $k_x = 4$. This makes integration order-dependent.
\begin{equation}
    \gamma_i = \sum_{j:\, x_j > x_i} \mathcal{D}_{k_x}\!\big(\alpha_j,\, x_j - x_i\big)
             = \sum_{j:\, x_j > x_i} \alpha_j \, 2^{-(x_j - x_i)/2^{k_x}}.
\end{equation}
A synaptic event at synapse $i$ carrying presynaptic burst $b$ (the number of APs) leaks the branch
voltage $V_B$ since its last event, then adds the gated EPSP ($k_B = 9$):
\begin{equation}
    V_B \leftarrow \mathcal{D}_{k_B}\!\big(V_B,\, \Delta t\big) \;+\; b\, w_i \,(1 + \gamma_i).
\end{equation}
The basal branch is a hard threshold: on crossing $\theta_B$ it emits a dendritic spike and resets.
\begin{equation}
    V_B \ge \theta_B \;\Longrightarrow\; \text{dendritic spike}, \quad V_B \leftarrow 0.
\end{equation}

\subsection{Apical dendrite}
The apical branch integrates $V_B$ identically (with $k = 11$) but \emph{does not} reset --- it leaks.
Instead of a discrete spike it delivers a graded plateau to the soma, a sigmoidal transfer
(Payeur et al., 2021) with half-activation $\theta_B$, ceiling $\delta V_S$ and slope
$\kappa = \ln 2 / 2^{k}$:
\begin{equation}
    \sigma^{\mathrm{ap}}(V_B) = \frac{\delta V_S}{1 + e^{-\kappa\,(V_B - \theta_B)}}, \qquad
    \sigma^{\mathrm{ap}} \in [0,\, \delta V_S].
\end{equation}

\subsection{Dendrite $\to$ soma coupling}
A basal dendritic spike with branch constant $c$ delivers $v_S = \max(c, 1)$ to the soma. Distal
branches ($c \le 0$) attenuate to unity at the soma but reinforce their own active synapses locally:
\begin{equation}
    v_S = \max(c,\,1), \qquad
    \alpha_j \leftarrow \alpha_j + |c| \quad (\text{for } \alpha_j > H_\alpha).
\end{equation}
An apical event delivers $v_S = \sigma^{\mathrm{ap}}(V_B)$.

\subsection{Soma}
The soma holds potential $V_S \in \mathfrak{i}8$ and a 6-bit burst counter $\beta \in [0, 63]$. An
incoming $v_S$ leaks the potential ($k_S = 10$) and the counter ($T_\beta = 500$ ticks per unit) over
$\Delta t$, then integrates:
\begin{equation}
    V_S \leftarrow \mathcal{D}_{k_S}\!\big(V_S,\, \Delta t\big) + v_S, \qquad
    \beta \leftarrow \beta - \left\lfloor \tfrac{\Delta t}{T_\beta} \right\rfloor.
\end{equation}
If the threshold $\theta_S > 0$ is crossed the soma emits a somatic spike whose burst $n$ counts how
many multiples of threshold accrued; it then resets to $V_{\mathrm{reset}} = -32$ and the burst
reinforces $\beta$:
\begin{equation}
    V_S \ge \theta_S \;\Longrightarrow\;
    n = \left\lfloor \tfrac{V_S}{\theta_S} \right\rfloor, \quad
    V_S \leftarrow V_{\mathrm{reset}}, \quad
    \beta \leftarrow \min(\beta + n,\; 63).
\end{equation}

\subsection{Plasticity (back-propagating AP)}
A somatic spike drives burst-dependent plasticity across the neuron's own afferent synapses. Only
sufficiently active synapses ($\alpha_i > H_\alpha$, with $H_\alpha = 30$) update; the burst counter
sets the sign about the baseline $H_\beta = 4$ ($\beta > H_\beta$: LTP; $\beta < H_\beta$: LTD), and
$\eta$ is the per-neuron learning rate ($\eta \ge 120$):
\begin{equation}
    \Delta w_i = \frac{(\beta - H_\beta)\,\alpha_i}{\eta}, \qquad
    w_i \leftarrow \mathrm{clamp}\big(w_i + \Delta w_i,\; -127,\; 127\big).
\end{equation}

\subsection{Event flow}
One spike threads a fixed sequence of events, the burst count $b$ riding the payload throughout:
\begin{equation}
    \texttt{SYNAPSE\_SIGNAL} \to
    \begin{cases}
        \texttt{DENDRITIC\_SPIKE} \to \texttt{SOMA\_SIGNAL} & \text{(basal)}\\
        \texttt{SOMA\_SIGNAL} & \text{(apical)}
    \end{cases}
    \to \texttt{SOMATIC\_SPIKE} \to \texttt{SYNAPSE\_SIGNAL}.
\end{equation}
A somatic spike fans out to every axonal target as an independent \texttt{SYNAPSE\_SIGNAL}.

\end{document}